\documentclass{beamer}
\usepackage[utf8]{inputenc}

\usetheme{Madrid}
\usecolortheme{default}
\usepackage{amsmath,amssymb,amsfonts,amsthm}
\usepackage{txfonts}
\usepackage{tkz-euclide}
\usepackage{listings}
\usepackage{adjustbox}
\usepackage{array}
\usepackage{tabularx}
\usepackage{gvv}
\usepackage{lmodern}
\usepackage{circuitikz}
\usepackage{tikz}
\usepackage{graphicx}
\usepackage{multicol}
\setbeamertemplate{page number in head/foot}[totalframenumber]

\usepackage{tcolorbox}
\tcbuselibrary{minted,breakable,xparse,skins}



\definecolor{bg}{gray}{0.95}
\DeclareTCBListing{mintedbox}{O{}m!O{}}{%
  breakable=true,
  listing engine=minted,
  listing only,
  minted language=#2,
  minted style=default,
  minted options={%
    linenos,
    gobble=0,
    breaklines=true,
    breakafter=,,
    fontsize=\small,
    numbersep=8pt,
    #1},
  boxsep=0pt,
  left skip=0pt,
  right skip=0pt,
  left=25pt,
  right=0pt,
  top=3pt,
  bottom=3pt,
  arc=5pt,
  leftrule=0pt,
  rightrule=0pt,
  bottomrule=2pt,

  colback=bg,
  colframe=orange!70,
  enhanced,
  overlay={%
    \begin{tcbclipinterior}
    \fill[orange!20!white] (frame.south west) rectangle ([xshift=20pt]frame.north west);
    \end{tcbclipinterior}},
  #3,
}
\lstset{
    language=C,
    basicstyle=\ttfamily\small,
    keywordstyle=\color{blue},
    stringstyle=\color{orange},
    commentstyle=\color{green!60!black},
    numbers=left,
    numberstyle=\tiny\color{gray},
    breaklines=true,
    showstringspaces=false,
}
%------------------------------------------------------------
%This block of code defines the information to appear in the
%Title page
\title %optional
{1.1.12}
\date{October  2025}
%\subtitle{A short story}

\author % (optional)
{BEERAM MADHURI - EE25BTECH11012}



\begin{document}


\frame{\titlepage}
\begin{frame}{Question}
Given closed-loop transfer function $T(s) = \frac{s-5}{(s+2)(s+3)}$, the system is: \hfill (EC 2007)

\begin{enumerate}
    \item[a)] Unstable
    \item[b)] Uncontrollable
    \item[c)] Minimum phase
    \item[d)] Non-minimum phase
\end{enumerate}
\end{frame}

\begin{frame}{solution}
    \frametitle{finding the Property of given system}
Given the closed-loop transfer function:
\begin{align}
    T(s) = \frac{s - 5}{(s + 2)(s + 3)}
\end{align}
Identifiying Poles and Zeros\\
\textbf{Poles:}\\
Set the denominator to zero: \\
\begin{align}
    (s + 2)(s + 3) = 0\\
    s = -2 ,-3
\end{align}
 Both are in the Left Half Plane (LHP).
\end{frame}
\begin{frame}
    \textbf{Zeros:}\\
Set the numerator to zero: 
\begin{align}{}
    s - 5 = 0\\
    s = +5
\end{align}
This is in the Right Half Plane (RHP).
\end{frame}
\begin{frame}{}
    Given Options:\\
a) Unstable:\\
A system is unstable if any of its poles lie in the Right Half Plane (RHP) or if there are repeated poles on the imaginary axis. (Our system is stable because poles are at -2 and -3).\\

Examples:\\
Stable: A commercial plane naturally stays level. If hit by wind, it corrects itself.\\
Unstable: A fighter jet without its computer would flip over and crash immediately.\\
\end{frame}
\begin{frame}{}
b)Uncontrollable:\\
A property related to State-Space analysis. In transfer functions, it usually implies a pole-zero cancellation. Since no terms cancel out here, we cannot assume it is uncontrollable from this form.\\

Examples:\\
Controllable: Every joint in a robotic arm moves when we tell it to. we can reach any point in space.\\

Uncontrollable: If a motor breaks, certain positions become impossible to reach, no matter what buttons we press.\\
\end{frame}
\begin{frame}
c)Minimum Phase:\\
A system where all poles and all zeros lie in the Left Half Plane (LHP). These systems have the minimum possible phase shift for a given magnitude response.\\

Examples:\\
Minimum Phase: You turn a car steering wheel right, and the car moves right immediately.\\
\end{frame}
\begin{frame}
d)Non-minimum Phase:\\
A system that has at least one zero (or pole) in the Right Half Plane (RHP). These systems exhibit "undershoot" in their step response and have more phase lag than a minimum phase system.\\

Examples:\\
Non-Minimum Phase: Like a large ship or forklift. To turn right, the back must swing left first. It moves the wrong way before going the right way.
\end{frame}
\begin{frame}
Determining the System Property:
Since we have a zero at $s = 5$, the system is non-minimum phase.\\
Correct Option: d) Non-minimum phase
\end{frame}

\begin{frame}[fragile]
\frametitle{C Code}
\begin{lstlisting}
#include <stdio.h>
#include <complex.h>
#include <math.h>

#define MAX_DEGREE 10
#define ITERATIONS 1000

// Function to evaluate polynomial at a complex point x
double complex evaluate(double coeff[], int degree, double complex x) {
    double complex res = coeff[degree];
    for (int i = degree - 1; i >= 0; i--) {
        res = res * x + coeff[i];
    }
    return res;
}
\end{lstlisting}
\end{frame}

\begin{frame}[fragile]
\frametitle{C Code}
\begin{lstlisting}
// Root finding using Durand-Kerner algorithm
void find_roots(double coeff[], int degree, double complex roots[]) {
    if (degree == 0) return;
    
    // Initial guesses for roots (complex numbers on a circle)
    for (int i = 0; i < degree; i++) {
        roots[i] = cpow(0.4 + 0.9 * I, i);
    }
\end{lstlisting}
\end{frame}

\begin{frame}[fragile]
\frametitle{C Code}
\begin{lstlisting}
    for (int iter = 0; iter < ITERATIONS; iter++) {
        for (int i = 0; i < degree; i++) {
            double complex denom = 1.0;
            for (int j = 0; j < degree; j++) {
                if (i != j) denom *= (roots[i] - roots[j]);
            }
            roots[i] -= evaluate(coeff, degree, roots[i]) / (coeff[degree] * denom);
        }
    }
}
\end{lstlisting}
\end{frame}

\begin{frame}[fragile]
\frametitle{C Code}
\begin{lstlisting}
int main() {
    int n_deg, d_deg;
    double num[MAX_DEGREE], den[MAX_DEGREE];
    double complex zeros[MAX_DEGREE], poles[MAX_DEGREE];

    printf("Enter degree of Numerator: ");
    scanf("%d", &n_deg);
    printf("Enter coefficients (from constant to highest power):\n");
    for(int i = 0; i <= n_deg; i++) scanf("%le", &num[i]);
\end{lstlisting}
\end{frame}

\begin{frame}[fragile]
\frametitle{C Code}
\begin{lstlisting}
    printf("\nEnter degree of Denominator: ");
    scanf("%d", &d_deg);
    printf("Enter coefficients (from constant to highest power):\n");
    for(int i = 0; i <= d_deg; i++) scanf("%le", &den[i]);

    find_roots(num, n_deg, zeros);
    find_roots(den, d_deg, poles);
\end{lstlisting}
\end{frame}

\begin{frame}[fragile]
\frametitle{C Code}
\begin{lstlisting}
    printf("TYPE,REAL,IMAG\n");
    for(int i = 0; i < n_deg; i++) 
        printf("ZERO,%.5f,%.5f\n", creal(zeros[i]), cimag(zeros[i]));
    for(int i = 0; i < d_deg; i++) 
        printf("POLE,%.5f,%.5f\n", creal(poles[i]), cimag(poles[i]));

    return 0;
}
\end{lstlisting}
\end{frame}

\begin{frame}[fragile]
    \frametitle{Python Code}
    \begin{lstlisting}
import sys
import pandas as pd
import matplotlib.pyplot as plt
import io

def plot_from_stdin():
    # Read the C output from standard input
    data = sys.stdin.read()
    df = pd.read_csv(io.StringIO(data))
    plt.figure(figsize=(8, 6))   
    # Separate poles and zeros
    zeros = df[df['TYPE'] == 'ZERO']
    poles = df[df['TYPE'] == 'POLE']
\end{lstlisting}
\end{frame}

\begin{frame}[fragile]
\frametitle{Python Code}
\begin{lstlisting}
    plt.scatter(zeros['REAL'], zeros['IMAG'], marker='o', s=100, facecolors='none', edgecolors='blue', label='Zeros')
    plt.scatter(poles['REAL'], poles['IMAG'], marker='x', s=100, color='red', label='Poles')

    # Formatting the S-Plane
    plt.axhline(0, color='black', lw=1) # Real axis
    plt.axvline(0, color='black', lw=1) # Imaginary axis
\end{lstlisting}
\end{frame}

\begin{frame}[fragile]
\frametitle{Python Code}
\begin{lstlisting}
    plt.grid(True, which='both', linestyle='--', alpha=0.5)
    plt.title('Pole-Zero Map (S-Plane)')
    plt.xlabel('Real (σ)')
    plt.ylabel('Imaginary (jω)')
    plt.legend()

    plt.axis('equal')
    plt.show()

if __name__ == "__main__":
    plot_from_stdin()
\end{lstlisting}
\end{frame}

\begin{frame}{Plot}
\begin{figure}[H]
\centering
\includegraphics[width=0.75\columnwidth]{figures/plot.png}
\caption{Poles and Zeroes}
\label{fig:placeholder}
\end{figure}
\end{frame}

\end{document}